\section*{Capítulo \ref{ch:b3:bioinformatics} --- Bioinformática y análisis de secuencias}

\begin{solution}{exo:b3:bioinformatics:1}
Global: las dos secuencias alineadas de extremo a extremo, con todo
residuo en una columna --- para dos proteínas que se creen homólogas en
toda su longitud, como los \omterm{def:b3:genomics:comparative}{ortólogos} de una enzima constitutiva. Local:
el par de subcadenas de mejor puntuación, ignorando el resto --- para
encontrar un dominio compartido (un dominio SH2 en dos proteínas de
señalización por lo demás no emparentadas), o un gen en una secuencia
genómica larga.
\end{solution}

\begin{solution}{exo:b3:bioinformatics:2}
Bordes $0,-1,-2,-3$ en ambos sentidos. Fila A: $1, 0, -1$.
Fila G: $0, 0, -1$. Fila C: $-1, -1, 1$. Óptimo $F(3,3) = 1$: AGC sobre
AAC sin huecos (coincidencia, discrepancia, coincidencia:
$1 - 1 + 1 = 1$).
\end{solution}

\begin{solution}{exo:b3:bioinformatics:3}
$s(a,a) = \lambda^{-1}\log\bigl(q_{aa}/p_{a}^{2}\bigr)$. El triptófano es
raro ($p_{W} \approx 0.013$), de modo que la probabilidad de que dos
triptófanos se alineen al azar, $p_{W}^{2}$, es minúscula, y un par de
triptófanos conservado es un indicio de homología mucho más fuerte que un
par de leucinas conservado ($p_{L}
\approx 0.1$); la razón de log-verosimilitudes es proporcionalmente
mayor.
\end{solution}

\begin{solution}{exo:b3:bioinformatics:4}
El \omterm{thm:b3:bioinformatics:evalue}{valor E} es el número de \omterm{def:b3:bioinformatics:alignment}{alineamientos} con una puntuación al menos tan
alta que cabría esperar por azar en una búsqueda de esta consulta contra
una base de datos de este tamaño. $E = 3$ significa que se esperan tres
puntuaciones así por azar: el acierto no es prueba de homología (puede
serlo de todos modos, pero la búsqueda no lo puede decir).
\end{solution}

\begin{solution}{exo:b3:bioinformatics:5}
$mn = 400\times 2\times 10^{11} = 8\times 10^{13} = 2^{46.2}$. $E(45) =
2^{1.2} \approx 2.3$; $E(55) = 2^{-8.8} \approx 2\times 10^{-3}$;
$E(65) = 2^{-18.8} \approx 2\times 10^{-6}$. $E = 10^{-3}$ exige $S' =
46.2 + 10.0 = 56$ bits. Una consulta de \num{40} residuos tiene un $mn$
diez veces menor, $2^{42.9}$: bastan $53$ bits --- pero una consulta
corta rara vez llega siquiera a eso.
\end{solution}

\begin{solution}{exo:b3:bioinformatics:6}
\omterm{def:b3:bioinformatics:motif}{Contenidos de información}: $2$, $1$, $0$, y $2 - H$ con $H = -(0.7\log_{2}
0.7 + 3\times 0.1\log_{2} 0.1) = 0.36 + 1.00 = 1.36$, o sea $0.64$. Total
$R = 3.64$ bits. Coincidencias por azar: $9.2\times 10^{6}$ posiciones en
dos hebras $\times 2^{-3.64} \approx 7\times 10^{5}$ --- el \omterm{def:b3:bioinformatics:motif}{motivo} es
casi inútil por sí solo.
\end{solution}

\begin{solution}{exo:b3:bioinformatics:7}
Una inserción de varios residuos es un solo suceso mutacional, de modo
que su coste no debería crecer linealmente con su longitud: un coste de
apertura más un pequeño coste de extensión lo modela. Una penalización de
apertura demasiado alta fuerza discrepancias donde corresponde un hueco y
desalinea todo lo que sigue a una inserción real; una demasiado baja
esparce huecos por todas partes, empareja residuos por azar e infla la
identidad.
\end{solution}

\begin{solution}{exo:b3:bioinformatics:8}
No necesariamente: el \omterm{def:b3:bioinformatics:alignment}{alineamiento} cubre un segmento de \num{130}
residuos, que es la firma de un dominio compartido más que la de un
\omterm{def:b3:genomics:comparative}{ortólogo} alineado en toda su longitud. Prueba: búsquese la proteína de la
mosca contra el proteoma humano (¿es la consulta su mejor acierto, en
toda la longitud?), identifíquese el dominio con un HMM de \omterm{def:b3:bioinformatics:msa}{perfil} y
constrúyase un árbol génico de la familia en varias especies.
\end{solution}

\begin{solution}{exo:b3:bioinformatics:9}
Un \omterm{def:b3:bioinformatics:msa}{perfil} registra, columna a columna, lo que la familia tolera: una
posición que es siempre hidrófoba pero nunca el mismo residuo, un residuo
catalítico invariante, una posición que es siempre un hueco en la mitad
de la familia. Un \omterm{def:b3:bioinformatics:alignment}{alineamiento} por pares puntúa cada residuo contra otro
único residuo y no puede saber que una valina en la posición 40 es ``tan
buena como'' la isoleucina que hay ahí en la consulta. El \omterm{def:b3:bioinformatics:msa}{perfil} además
pondera las columnas conservadas, de modo que una similitud débil
concentrada allí donde la familia está conservada se vuelve significativa.
\end{solution}

\begin{solution}{exo:b3:bioinformatics:10}
Con una \omterm{prop:b3:bioinformatics:logodds}{puntuación esperada} positiva $\mu > 0$ por columna, la puntuación
acumulada a lo largo de la diagonal de dos secuencias aleatorias es un
paseo aleatorio con deriva positiva: tras $n$ columnas vale unos $\mu n$,
de modo que el mejor \omterm{def:b3:bioinformatics:alignment}{alineamiento local} es esencialmente el todo y su
puntuación crece como $\mu n$ y no como $\log n$. La teoría de
Karlin--Altschul, que exige una deriva negativa para que las
puntuaciones altas sean excursiones raras, no se aplica y no existe
ningún $\lambda$. Para ADN al \qty{60}{\%} de GC la probabilidad de una
coincidencia es $2(0.3^{2}) + 2(0.2^{2}) = 0.26$, de modo que la
\omterm{prop:b3:bioinformatics:logodds}{puntuación esperada} es $0.26 - 0.74 = -0.48$: sigue siendo negativa y la
estadística vale; pero un esquema como coincidencia $+1$ y discrepancia
$-0.3$ tendría esperanza $+0.04$ y daría todo el genoma como un solo
\omterm{def:b3:bioinformatics:alignment}{alineamiento}.
\end{solution}

\begin{solution}{exo:b3:bioinformatics:11}
\omterm{def:b3:bioinformatics:blast}{Semillas} y encadenamiento: se buscan coincidencias exactas o casi exactas
de $k$-meros entre las dos secuencias con una tabla de dispersión, se
conservan las que se alinean en diagonales coherentes, se encadenan y se
aplica \omterm{thm:b3:bioinformatics:nw}{programación dinámica} solo en los huecos entre \omterm{def:b3:bioinformatics:blast}{semillas}
encadenadas; se renuncia a los \omterm{def:b3:bioinformatics:alignment}{alineamientos} en regiones sin \omterm{def:b3:bioinformatics:blast}{semilla}
(tramos muy divergentes). Bandas: si se sabe que las dos secuencias son
casi colineales, se calculan solo las celdas de una banda de anchura $w$
alrededor de la diagonal, a un coste $wn$ en lugar de $mn$; se renuncia a
cualquier \omterm{def:b3:bioinformatics:alignment}{alineamiento} con una inserción mayor que la banda.
\end{solution}

\begin{solution}{exo:b3:bioinformatics:12}
La secuencia codificante tiene un período de tres: las tres posiciones
del codón tienen composiciones de bases distintas (la tercera es la más
sesgada), y el uso de codones es desigual en cada especie. Un modelo con
tres estados codificantes en serie, cada uno emitiendo bases con la
composición de esa posición del codón, asigna al ADN codificante una
probabilidad mayor que el estado no codificante, a lo largo de una
ventana de unas pocas decenas de codones, incluso sin paradas. En un
genoma humano los exones son cortos (\qty{150}{bp}) y están separados por
intrones de kilobases, de modo que la señal codificante es breve e
interrumpida; el modelo debe reconocer además los sitios de corte, que
son señales débiles, y la enorme cantidad de secuencia no codificante
produce muchos segmentos codificantes falsos.
\end{solution}

\begin{solution}{pb:b3:bioinformatics:1}
\textbf{1.} Filas K, Q, T; columnas K, A, Q, T; bordes $0,-1,-2,-3,-4$ y
$0,-1,-2,-3$. Fila K: $1, 0, -1, -2$; fila Q: $0, 0, 1, 0$; fila T:
$-1, -1, 0, 2$. Óptimo $2$: \texttt{K-QT} sobre \texttt{KAQT}.
\textbf{2.} Mejor puntuación local $3$: \texttt{CAT} contra \texttt{CAT}
(residuos 4--6 de GATCAT con 2--4 de ACAT); \texttt{ATCAT} contra
\texttt{A-CAT} también puntúa $4 - 1 = 3$.
\textbf{3.} $300\times 450 = 1.35\times 10^{5}$ actualizaciones; contra
la base de datos, $300\times 1.2\times 10^{11} = 3.6\times 10^{13}$.
\textbf{4.} $3.6\times 10^{4}$ s, diez horas por consulta; las \omterm{def:b3:bioinformatics:blast}{semillas}
de \omterm{def:b3:bioinformatics:blast}{BLAST} se saltan casi toda la tabla y responden en segundos.
\textbf{5.} $q_{ab} = p_{a}p_{b}e^{\lambda s}$. Alanina: $e^{1.39} = 4.0$,
$q_{AA} = 0.074^{2}\times 4.0 = 0.022$, razón $4$. Triptófano:
$e^{3.82} = 45$, $q_{WW} = 0.013^{2}\times 45 = 0.0077$, razón $45$. Un
par de triptófanos alineado es $45$ veces más frecuente en homólogos que
por azar, y un par de alaninas solo cuatro veces; aun así los pares de
alaninas son más frecuentes en términos absolutos porque la alanina es
frecuente.
\textbf{6.} El \qty{24}{\%} cae en la zona crepuscular, donde los
\omterm{def:b3:bioinformatics:alignment}{alineamientos} aleatorios alcanzan el \qtyrange{15}{20}{\%}; lo zanjarían
el \omterm{thm:b3:bioinformatics:evalue}{valor E} del \omterm{def:b3:bioinformatics:alignment}{alineamiento}, unos \omterm{def:b3:bioinformatics:motif}{motivos} conservados en las posiciones
correctas, una coincidencia con un HMM de \omterm{def:b3:bioinformatics:msa}{perfil} de una familia conocida
o un plegamiento compartido.
\textbf{7.} $mn = 300\times 1.2\times 10^{11} = 3.6\times 10^{13}$;
$\log_{2}(mn) = 45.0$.
\textbf{8.} $E(92) = 2^{45 - 92} = 2^{-47} \approx 7\times 10^{-15}$;
$E(38) = 2^{7} = 128$.
\textbf{9.} $E = 10^{-3}$ con $S' = 45 + 10 = 55$ bits; $E = 1$ con $45$
bits.
\textbf{10.} Diez veces $mn$: $E \approx 7\times 10^{-14}$, aún
abrumador.
\textbf{11.} Con $E = 128$, el décimo acierto es lo que produce el azar;
un \qty{40}{\%} de identidad a lo largo de $60$ residuos no es prueba.
Una búsqueda por \omterm{def:b3:bioinformatics:msa}{perfil} de los residuos 200--260 contra la base de datos
de dominios podría mostrar si ese segmento es un dominio conocido, con
una estadística de la que carece la comparación por pares.
\textbf{12.} Los mejores aciertos recíprocos en toda la longitud son
compatibles con una ortología uno a uno; no la demuestran --- una
duplicación en un linaje posterior a la separación da dos coortólogos, y
la pérdida del \omterm{def:b3:genomics:comparative}{ortólogo} verdadero puede dejar un \omterm{def:b3:genomics:comparative}{parálogo} como mejor
acierto. La prueba es un árbol génico con varias especies.
\textbf{13.} $R = 2 + 2 + 1.6 + 2 + 0.8 + 1.2 + 0.4 + 0.3 = 10.3$ bits.
\textbf{14.} $3.2\times 10^{9}\times 2^{-10.3} \approx 2.5\times 10^{6}$
coincidencias por azar.
\textbf{15.} $\log_{2}(3.2\times 10^{9}/200) = \log_{2}(1.6\times 10^{7})
\approx 24$ bits.
\textbf{16.} La coaparición con un espaciado fijo suma los bits: $10.3 + 8 =
18.3$, lo que aporta $8$ de los $13.7$ que faltan; unos $5.7$ bits (un
factor de $50$ en las coincidencias por azar) han de venir de otro sitio
--- la accesibilidad de la \omterm{def:b3:chromatin-epigenetics:nucleosome}{cromatina}, más compañeros.
\textbf{17.} $H = -(0.5\log_{2}0.5 + 0.5\log_{2}0.5) = 1$ bit; $R = 2 -
1 = 1$ bit.
\textbf{18.} Un factor bacteriano ha de encontrar sus pocos sitios en un
genoma de \qty{4.6}{Mb} sin ayuda de la \omterm{def:b3:chromatin-epigenetics:nucleosome}{cromatina}: necesita unos $19$
bits, y sus sitios son largos y conservados. Un genoma eucariota es mil
veces mayor y exige diez bits más, y sin embargo sus factores tienen
sitios cortos; consiguen la especificidad por combinación y por la
restricción a la \omterm{def:b3:chromatin-epigenetics:nucleosome}{cromatina} accesible, lo que además hace la regulación
más evolucionable, ya que un sitio corto se gana o se pierde con
facilidad.
\textbf{19.} $3/64 = 0.047$; el número esperado de codones antes de una
parada es $64/3 \approx 21$.
\textbf{20.} $p_{A} = p_{T} = 0.31$, $p_{G} = p_{C} = 0.19$: $P(\text{TAA})
= 0.31^{3} = 0.030$, $P(\text{TAG}) = P(\text{TGA}) = 0.31^{2}\times 0.19
= 0.018$; total $0.066$, longitud esperada del marco, $15$ codones. El
ADN rico en AT está lleno de paradas, de modo que los marcos abiertos al
azar son más cortos y los largos destacan más.
\textbf{21.} $(61/64)^{100} = e^{-4.80} = 0.008$; $(61/64)^{300} =
e^{-14.4} = 5.6\times 10^{-7}$.
\textbf{22.} Cada marco abierto maximal termina en una parada, y $3.2\times
10^{9}$ inicios de codón contienen $3.2\times 10^{9}\times 3/64 = 1.5\times
10^{8}$ paradas: unos $1.5\times 10^{8}\times 0.008 = 1.2\times 10^{6}$
marcos al azar de al menos $100$ codones, y $1.5\times 10^{8}\times
5.6\times 10^{-7} \approx 80$ de al menos $300$.
\textbf{23.} Una bacteria de \qty{4.6}{Mb} tiene unas $4\times 10^{5}$
paradas y por tanto unos \num{3500} marcos por azar de $100$ codones pero
casi ninguno de $300$; sus genes promedian $300$ codones y el
\qty{88}{\%} del ADN es codificante, de modo que un marco abierto largo
es casi siempre un gen. En el gusano codifica el \qty{1.5}{\%} del ADN,
los exones promedian $50$ codones --- menos que el umbral del azar --- y
un millón de marcos al azar de $100$ codones los sepultan. Los buscadores
de genes eucariotas usan señales de sitios de corte, el sesgo de codones
en un modelo oculto de Markov, la homología con proteínas conocidas y,
sobre todo, los transcritos secuenciados.
\textbf{24.} Una lectura de un mensajero cortado y empalmado se alinea
con el genoma en dos trozos separados por un intrón: la partición marca
los dos sitios de corte hasta la base; la cobertura de lecturas delimita
los exones y las lecturas emparejadas enlazan exones sucesivos en un solo
transcrito, lo que resuelve cuál de varios sitios de corte candidatos se
usa.
\textbf{25.} $E \approx 7\times 10^{-15}$ para el mejor acierto; unos
$24$ bits para especificar $200$ sitios en el genoma; y unos $80$ marcos
de lectura por azar de $300$ codones en todo el genoma.
\end{solution}
